\section{Institutions}\label{sec:inst}

\subsection{The Danish public pension}

Every Danish resident above the pension age receives the \emph{folkepension}. The pension age was 65 until mid-2019 and has risen in steps since \citep{DKSTARalder}; for simplicity, and because the change affects only the last wave and only a few months of it, older people are defined here as those aged 65 and over. The folkepension has two parts. The basic amount, \emph{grundbeløb}, is the same for everyone apart from a reduction for high earnings from work. The pension supplement, \emph{pensionstillæg}, is paid at a higher rate to people who are single in the sense of the pension law than to those who live in a couple, and is withdrawn as income from sources other than the folkepension rises. A single pensioner with little other income therefore receives a fixed annual sum, the full basic amount plus the full supplement for singles; these sums are given for the LIS years in Table~\ref{tab:floors}. Both parts are taxable income. Since 2003 pensioners with little liquid wealth and little other income have also received a supplementary payment, the \emph{supplerende pensionsydelse} or \emph{ældrecheck}, a flat annual sum (10{,}700 kroner in 2010) that is also taxable \citep{DKVejledning2010}. Occupational and private pensions and the statutory funded supplement ATP come on top of the folkepension and reduce the supplement.

\subsection{The Danish housing benefit for pensioners}

Pensioners who rent their home can receive \emph{boligydelse}, the pensioners' version of the Danish housing benefit \citep{DKBoligstotte2025}. It is computed from the rent, the household's income and wealth, the number of persons and the size of the dwelling: it covers a fixed share of the housing expense above a threshold, less a share of household income above an income threshold, subject to a minimum payment by the household and an annual maximum. Owners and members of housing cooperatives can receive it only as a loan repayable when the dwelling is sold. The benefit is not taxable. For people who are not pensioners the corresponding scheme, \emph{boligsikring}, is considerably less generous. The annual maximum for a pensioner was 45{,}480 kroner in 2016 and 48{,}540 kroner in 2019 \citep{DKVejledning2016,DKVejledning2019}.

Two features of the benefit matter here. It is received by some pensioners and not by others according to a fact, renting rather than owning, that is only loosely related to the rest of their income. And because it is tax-free while the pension is taxed, it enters disposable income in full. A single pensioner who lives on the folkepension and rents therefore has a disposable income equal to the net pension plus a benefit of a few tens of thousands of kroner, while one who owns her home has the net pension alone.

\subsection{The Danish student grant}

Students in higher education receive a monthly grant from the state educational support scheme, SU, at a higher rate if they do not live with their parents. The rates for a student living away from home are given in Table~\ref{tab:floors} as twelve times the monthly rate. The grant is taxable \citep{DKSkatSU} and exceeds the personal allowance, so a student who lives on it has a disposable income somewhat below the grant.

\subsection{The Norwegian minimum pension}

The Norwegian national insurance scheme pays an old-age pension, drawn by most people at 67, that depends on past earnings. A person whose earnings-related entitlement is low receives a top-up to a guaranteed minimum, the \emph{minste pensjonsnivå}, which is adjusted every year and has had a separate, higher rate for single pensioners since September 2016 \citep{NOMeld2016,NOMeld2020,NOMeld2022}. Table~\ref{tab:floors} gives calendar-year averages of the rate for a single pensioner. The Norwegian tax rules for pensioners are designed so that a person whose only income is the minimum pension pays no income tax \citep{NOSkatteetaten}; the disposable income of a minimum pensioner is therefore the minimum pension itself. Housing allowances exist in Norway too, through the means-tested \emph{bostøtte}, but the minimum pension is paid to far fewer people than the Danish full folkepension, because the earnings-related pension exceeds the minimum for almost everyone with a working history.

\begin{table}[t]
\centering\small
\caption{Statutory income floors for a single person, annual amounts in national currency (gross). Norway: calendar-year average of the minimum pension level for a single old-age pensioner (minste pensjonsnivå, særskilt sats from September 2016), computed from the rates in force and their effective dates. Denmark: folkepension basic amount plus the maximum pension supplement for a single pensioner; SU grant for a student living away from home in higher education, twelve months. Sweden: guarantee pension for an unmarried pensioner, $2.13$ price base amounts. Finland: full national pension plus guarantee pension top-up (from 2011) for a single person living alone, twelve months. Sources: \citet{DKTIMES} and the annual rate guidance \citep{DKVejledning2010} for the folkepension; \citet{DKSUsatser} for SU; \citet{NOMeld2016,NOMeld2020,NOMeld2022} for Norway; \citet{SEPensionsmyndigheten} and \citet{SCBpbb} for Sweden; \citet{ETK2024} for Finland.}\label{tab:floors}
\begin{tabular}{lrrrrrr}
\toprule
 & 2004 & 2007 & 2010 & 2013 & 2016 & 2019\\
\midrule
Norway, minimum pension (NOK) & 104{,}260 & 117{,}476 & 148{,}704 & 166{,}180 & 179{,}745 & 198{,}609\\
Denmark, folkepension (DKK) & 111{,}924 & 119{,}244 & 133{,}272 & 141{,}984 & 149{,}544 & 159{,}000\\
Denmark, SU grant (DKK) & 54{,}228 & 58{,}224 & 64{,}608 & 69{,}036 & 71{,}292 & 73{,}992\\
Sweden, guarantee pension (SEK) & 83{,}709 & 85{,}839 & 90{,}312 & 94{,}785 & 94{,}359 & 99{,}045\\
Finland, minimum pension (EUR) & 5{,}957 & 6{,}298 & 7{,}010 & 8{,}866 & 9{,}202 & 9{,}414\\
\bottomrule
\end{tabular}
\end{table}

